\chapter{Vollständigkeit und abgeschlossene Teilräume}
\label{b47guide:ch:completeness}

\section{Vollständigkeit erkennen}

\begin{definition}[Vollständigkeit]
\label{b47guide:def:complete}
$(X,d)$ heißt \emph{vollständig}, wenn jede Cauchy-Folge in $X$ einen
Grenzwert in $X$ besitzt. Eine Teilmenge $A\subseteq X$ heißt vollständig,
wenn $(A,d_A)$ vollständig ist. Der leere Raum ist vollständig.
\end{definition}

\begin{theorem}[Vollständigkeit über Teilfolgen]
Ein metrischer Raum ist genau dann vollständig, wenn jede seiner
Cauchy-Folgen eine im Raum konvergente Teilfolge besitzt.
\end{theorem}

\begin{theorem}[Vollständigkeit und Abgeschlossenheit]
\label{b47guide:thm:closed-complete}
Jeder vollständige Teilraum eines metrischen Raums ist abgeschlossen.
Ist $X$ vollständig, so gilt für jede Teilmenge $A\subseteq X$:
\[
 (A,d_A)\text{ vollständig}\quad\Longleftrightarrow\quad
 A\text{ abgeschlossen in }X.
\]
Insbesondere sind abgeschlossene Kugeln eines vollständigen Raums
vollständig. \cite{project,metric-notes}
\end{theorem}

\begin{theorem}[Cantorsches Vollständigkeitskriterium]
\label{b47guide:thm:cantor}
Für einen metrischen Raum $X$ sind äquivalent:
\begin{enumerate}
\item $X$ ist vollständig.
\item Für jede Folge nichtleerer abgeschlossener Teilmengen
$F_1\supseteq F_2\supseteq\cdots$ mit
$\operatorname{diam}F_n\to0$ besteht $\bigcap_{n\geq1}F_n$ aus genau einem
Punkt.
\end{enumerate}
Die Durchmesser dürfen zunächst unendlich sein; die Grenzbedingung
verlangt, dass sie schließlich beliebig klein und endlich sind.
\cite{baire-notes}
\end{theorem}

\section{Diskrete Abstände und Grenzen der Folgerungen}

\begin{definition}[Standarddiskrete Metrik und gleichmäßige Diskretheit]
Die \emph{standarddiskrete Metrik} auf $X$ ist
\[
 \delta(x,y)=\begin{cases}0,&x=y,\\1,&x\ne y.\end{cases}
\]
Ein metrischer Raum heißt \emph{gleichmäßig diskret}, wenn ein $c>0$
existiert mit $d(x,y)\geq c$ für alle verschiedenen $x,y\in X$.
\end{definition}

\begin{theorem}[Folgen und Vollständigkeit bei getrennten Punkten]
In einem gleichmäßig diskreten metrischen Raum sind die Cauchy-Folgen
genau die schließlich konstanten Folgen. Konvergenz gegen $x$ bedeutet,
dass die Folge schließlich gleich $x$ ist. Jeder solche Raum ist
vollständig; er ist genau dann total beschränkt beziehungsweise kompakt,
wenn er endlich ist.
\end{theorem}

\begin{theorem}[Die Standarddiskretmetrik trennt Beschränktheit und Kompaktheit]
Jeder standarddiskrete Raum ist beschränkt und vollständig. Ist sein
Träger unendlich, so ist er weder total beschränkt noch kompakt.
Insbesondere reicht auch die Kombination aus Vollständigkeit,
Abgeschlossenheit und Beschränktheit nicht als allgemeines metrisches
Kompaktheitskriterium. \cite{project}
\end{theorem}
\section{Metrische Produkte}

\begin{definition}[Endliche und abzählbare Produktmetrik]
Für metrische Räume $(X_j,d_j)$ ist auf $\prod_{j=1}^{m}X_j$, $m\geq1$,
die \emph{Maximumsmetrik} definiert durch
\[
 D_{\max}(x,y)=\max_{1\leq j\leq m}d_j(x_j,y_j).
\]
Auf $\prod_{j=1}^{\infty}X_j$ wird eine Metrik definiert durch
\[
 D(x,y)=\sum_{j=1}^{\infty}2^{-j}\min\{1,d_j(x_j,y_j)\}.
\]
\end{definition}

\begin{theorem}[Metrische Produkteigenschaften]
\label{b47guide:thm:products}
Die angegebenen Produktabstände sind Metriken. Konvergenz und die
Cauchy-Eigenschaft einer Folge bezüglich dieser Metriken sind jeweils
äquivalent zur entsprechenden Eigenschaft in jeder Koordinate.
Sind alle Faktoren vollständig, total beschränkt beziehungsweise kompakt,
so besitzt das Produkt die entsprechende Eigenschaft. Für nichtleere
Faktoren gelten auch die Umkehrungen. \cite{product-notes}
\end{theorem}